\section{Proposed Scheme}\label{sec:method}
\subsection{Problem Formulation}
Let $p_i=(\phi_i,\lambda_i,t_i)$ be a phone location and timestamp. A quality indicator $I_i$ marks a valid consecutive segment, and $d_g$ denotes a geographic distance. Observed trajectory distance and segment-average speed are
\begin{align}
D &= \sum_{i=2}^{n} I_i d_g(p_{i-1},p_i),\label{eq:distance}\\
v_i &= \frac{d_g(p_{i-1},p_i)}{t_i-t_{i-1}},\quad t_i>t_{i-1}.\label{eq:speed}
\end{align}
Missing intervals are marked and not treated as observed paths. The average speed must specify whether its denominator includes stationary time. If a reliable step count $S$ is available for interval $\Delta t$, cadence is $S/\Delta t$ with consistent units. Thresholds and smoothing parameters require experimental validation.

For a location $p$ and POI $q$, the direct distance $d_g(p,q)$ and route-network distance $d_{route}(p,q)$ are distinct. Route distance and accessibility are reported only when available from a route service. Approximate total energy expenditure in kcal is
\begin{equation}
E=\frac{MET\times3.5\times m_{kg}\times t_{min}}{200},\label{eq:energy}
\end{equation}
where activity-specific $MET$, body mass and duration require stated sources and assumptions. The estimate is not a physiological measurement. For food item $i$, user-confirmed quantity $g_i$ in grams and a database value $v_{ik}$ per 100 g yield nutrient $N_k=\sum_i(g_i/100)v_{ik}$. No heart rate or cardiac workload is inferred from these quantities.

\subsection{System Architecture and Data Flow}
The user opens a mobile web interface and states an exercise goal to an AI avatar. The avatar coordinates GIS queries for routes and nearby POIs; the system presents sources, distances and route limitations for user confirmation. During activity, the phone records permitted location and available motion data, applies quality checks and displays activity estimates. GIS maps and POIs are external data sources, whereas the phone is the required personal sensing device. Fig.~\ref{fig:architecture} shows the proposed user journey.

After activity, the intended design encrypts consented activity data for cloud storage. A digest is submitted to an anchoring interface only with separate consent; raw personal data are not placed on chain. The documented sensor prototype uses AES-256-GCM and a SHA-256 digest of the complete encrypted JSON package \cite{blockchain}. It currently disables real chain writes, and integration with GPS activity data and the complete coaching system is unverified. The user may also query food via MCP, confirm item and portion, and create a dietary record; its storage and authorization need implementation and validation. The food path is not assumed to share the encrypted sensor package.

\begin{figure*}[t]
\centering
\resizebox{0.96\textwidth}{!}{%
\begin{tikzpicture}[block/.style={draw,rounded corners,fill=blue!6,align=center,text width=29mm,minimum height=12mm,font=\small},arrow/.style={-{Stealth[length=2mm]},thick}]
\node[block] (web) at (0,0) {Mobile web\\and AI avatar};
\node[block] (route) at (4.2,0) {GIS route and\\POI planning};
\node[block] (confirm) at (8.4,0) {User confirms\\route and consent};
\node[block] (sensor) at (12.6,0) {Phone sensing\\during activity};
\node[block] (analysis) at (12.6,-2.6) {Trajectory and\\activity estimates};
\node[block] (cloud) at (8.4,-2.6) {Consented encrypted\\cloud storage};
\node[block] (chain) at (4.2,-2.6) {Separate consent:\\digest anchoring};
\node[block] (food) at (0,-2.6) {Food MCP query\\and confirmed record};
\draw[arrow] (web) -- (route);
\draw[arrow] (route) -- (confirm);
\draw[arrow] (confirm) -- (sensor);
\draw[arrow] (sensor) -- (analysis);
\draw[arrow] (analysis) -- (cloud);
\draw[arrow] (cloud) -- (chain);
\draw[arrow] (web) -- (food);
\end{tikzpicture}%
}
\caption{Proposed mobile-web user journey and data paths. Food recording is a separate path; digest anchoring requires separate consent and real chain writes are currently disabled.}
\label{fig:architecture}
\end{figure*}

A separate data-management flow (Fig.~\ref{fig:data-flow}) distinguishes encrypted sensor recordings from the separately consented public-summary path described in the prototype. Digest comparison requires a trusted reference, and the AES-GCM tag is verified during decryption; neither proves sensor truth. Local plaintext samples and recovery keys also require protection.
\begin{figure*}[t]
\centering
\resizebox{0.96\textwidth}{!}{%
\begin{tikzpicture}[block/.style={draw,rounded corners,fill=blue!6,align=center,text width=30mm,minimum height=12mm,font=\small},arrow/.style={-{Stealth[length=2mm]},thick}]
\node[block] (sample) at (0,0) {Phone samples\\and local storage};
\node[block] (encrypt) at (4.3,0) {Encrypt package\\and hash};
\node[block] (api) at (8.6,0) {Authenticated\\cloud API};
\node[block] (store) at (12.9,0) {Off-chain\\ciphertext};
\node[block] (anchor) at (8.6,-2.4) {Consent and\\anchoring interface};
\node[block] (public) at (0,-2.4) {Separately consented\\public batches};
\node[block] (summary) at (4.3,-2.4) {Public summary\\and validation};
\node[block] (assistant) at (4.3,-4.4) {AI assistant};
\draw[arrow] (sample) -- (encrypt);
\draw[arrow] (encrypt) -- (api);
\draw[arrow] (api) -- (store);
\draw[arrow] (api) -- (anchor);
\draw[arrow] (public) -- (summary);
\draw[arrow] (summary) -- (assistant);
\end{tikzpicture}%
}
\caption{Prototype data-management paths. Public batches are separate from encrypted recordings; actual chain transactions and complete integration are not established.}
\label{fig:data-flow}
\end{figure*}

\subsection{Proposed Algorithm}
\begin{algorithm}[t]
\caption{Mobile Exercise and Food Workflow}
\begin{algorithmic}[1]
\Require User request, permissions, available phone and tool data
\Ensure Route, activity summary and optional confirmed food record
\State Query GIS for candidate routes and POIs
\State Present route and obtain user confirmation
\State Collect permitted phone samples during activity
\State Check quality and missing intervals; calculate \eqref{eq:distance}--\eqref{eq:speed}
\State Estimate cadence and energy only if required inputs are available
\State Show nearby POIs, distance type and data limitations
\State With consent, encrypt and upload the activity package
\State Request digest anchoring only with separate consent
\If{user requests food information}
\State Query food MCP; confirm item and portion before recording
\EndIf
\State Report unavailable data, upload failures and anchoring status
\end{algorithmic}
\end{algorithm}
